%=============================================================================
\section{Protocol Stack Architecture}
%=============================================================================

\subsection{Layer Overview}

The STAR protocol implements a four-layer stack over SPI peripheral mode:

\begin{figure}[h]
\centering
\begin{tikzpicture}[
    layer/.style={draw, rectangle, minimum width=10cm, minimum height=1.2cm, font=\small},
    node distance=0cm
]
    \node[layer, fill=red!20] (app) {
        \textbf{Application Layer} - Motor Commands, Telemetry, OTA
    };
    \node[layer, fill=orange!20, below=of app] (serial) {
        \textbf{Serialization Layer} - nanopb (Protocol Buffers)
    };
    \node[layer, fill=yellow!20, below=of serial] (arq) {
        \textbf{ARQ Layer} - Stop-and-Wait (SEQ, ACK/NACK, Timeout 8ms)
    };
    \node[layer, fill=green!20, below=of arq] (frame) {
        \textbf{Frame Layer} - SYNC + Header (8B) + Payload + CRC-32 (4B)
    };
    \node[layer, fill=blue!20, below=of frame] (transport) {
        \textbf{Transport Layer} - SPI Peripheral (10~Mbps) + DMA
    };
\end{tikzpicture}
\caption{STAR Protocol Stack Layers}
\end{figure}

\subsection{Design Requirements}

\begin{table}[h]
\centering
\caption{Protocol Performance Targets}
\begin{tabular}{lcc}
\toprule
\textbf{Metric} & \textbf{Target} & \textbf{Rationale} \\
\midrule
Memory footprint & <10~KB & 2\% of 512~KB SRAM \\
Reliability & 99.9999\% & CRC-32 error detection \\
Throughput margin & 27× & Stop-and-Wait sufficient \\
Emergency stop latency & <20~ms & Safety requirement \\
Dynamic allocation & Zero & Safety-critical system \\
SPI speed & 10~Mbps & Custom PCB <15cm traces \\
\bottomrule
\end{tabular}
\end{table}

%=============================================================================
\section{Phase 1: CRC Library (lib/star\_crc)}
%=============================================================================

\subsection{Hardware CRC-32 Acceleration}

The ESP32-S3 provides ROM-based CRC-32 calculation using \texttt{esp\_rom\_crc32\_le()} with IEEE 802.3 polynomial (0x04C11DB7). This is 5-10× faster than software CRC-16.

\subsubsection{Directory Structure}

\begin{verbatim}
lib/star_crc/
├── include/
│   └── star_crc.h
├── src/
│   └── star_crc.c
├── library.json
└── CMakeLists.txt
\end{verbatim}

\subsubsection{API Definition}

\begin{lstlisting}[caption={star\_crc.h - CRC-32 Interface}]
#ifndef STAR_CRC_H
#define STAR_CRC_H

#include <stdint.h>
#include <stddef.h>
#include "esp_err.h"

/**
 * @brief Initialize CRC-32 module
 * @return ESP_OK on success
 */
esp_err_t star_crc_init(void);

/**
 * @brief Calculate CRC-32 checksum
 * @param data Pointer to data buffer
 * @param len Length of data in bytes
 * @return CRC-32 checksum (IEEE 802.3)
 */
uint32_t star_crc_calculate(const uint8_t* data, size_t len);

/**
 * @brief Validate CRC-32 checksum
 * @param data Pointer to data buffer
 * @param len Length of data in bytes
 * @param expected_crc Expected CRC-32 value
 * @return true if CRC matches, false otherwise
 */
bool star_crc_validate(const uint8_t* data, size_t len,
                       uint32_t expected_crc);

#endif // STAR_CRC_H
\end{lstlisting}

\subsubsection{Implementation}

\begin{lstlisting}[caption={star\_crc.c - CRC-32 Implementation}]
#include "star_crc.h"
#include "rom/crc.h"  // ESP32-S3 ROM CRC functions
#include "esp_log.h"

static const char* TAG = "star_crc";

esp_err_t star_crc_init(void) {
    ESP_LOGI(TAG, "CRC-32 module initialized (hardware-accelerated)");
    return ESP_OK;
}

uint32_t star_crc_calculate(const uint8_t* data, size_t len) {
    if (data == NULL || len == 0) {
        return 0;
    }

    // Use ESP32-S3 ROM CRC-32 (IEEE 802.3 polynomial)
    return crc32_le(0, data, len);
}

bool star_crc_validate(const uint8_t* data, size_t len,
                       uint32_t expected_crc) {
    uint32_t calculated_crc = star_crc_calculate(data, len);
    return (calculated_crc == expected_crc);
}
\end{lstlisting}

%=============================================================================
\section{Phase 2: Protocol Buffers (lib/star\_proto)}
%=============================================================================

\subsection{Directory Structure}

\begin{verbatim}
lib/star_proto/
├── include/
│   ├── star_proto.h        # Main protocol interface
│   ├── star_frame.h        # Frame layer API
│   └── star_arq.h          # ARQ layer API
├── src/
│   ├── star_frame.c        # Frame implementation
│   ├── star_arq.c          # ARQ implementation
│   └── generated/          # nanopb-generated code
│       ├── messages.pb.h
│       └── messages.pb.c
├── proto/
│   ├── messages.proto      # Protobuf message schemas
│   └── messages.options    # nanopb configuration
├── library.json
└── CMakeLists.txt
\end{verbatim}

\subsection{Message Schemas (messages.proto)}

\begin{lstlisting}[caption={messages.proto - Protocol Buffer Definitions},language=Python]
syntax = "proto3";

// ============ Commands (RPi5 -> ESP32) ============

message MotorSpeedCmd {
    uint32 motor_id = 1;        // 0-3
    float setpoint_rpm = 2;     // Target RPM
}

message MotorStopCmd {
    uint32 motor_id = 1;        // 0-3
    bool use_brake = 2;         // true=brake, false=coast
}

message EmergencyStopCmd {
    // No fields - immediate stop all motors
}

message SetPidGainsCmd {
    uint32 motor_id = 1;        // 0-3
    float kp = 2;
    float ki = 3;
    float kd = 4;
}

message OtaWriteChunkCmd {
    uint32 offset = 1;          // Byte offset in firmware
    bytes data = 2;             // Chunk data (max 1024 bytes)
}

// ============ Responses (ESP32 -> RPi5) ============

message TelemetryResp {
    repeated float motor_rpm = 1;        // 4 motors
    repeated float motor_current_ma = 2; // 4 motors
    float temperature_c = 3;
    float voltage_v = 4;
}

message StatusResp {
    uint32 state = 1;           // System state
    uint32 error_code = 2;      // Error code (0 = no error)
}

message AckResp {
    uint32 seq_number = 1;      // Acknowledged sequence number
}

message NackResp {
    uint32 seq_number = 1;      // Failed sequence number
    uint32 error_code = 2;      // Reason for NACK
}
\end{lstlisting}

\subsection{Nanopb Options (messages.options)}

\begin{lstlisting}[caption={messages.options - Static Allocation Config}]
# Maximum array sizes (static allocation)
TelemetryResp.motor_rpm max_count:4
TelemetryResp.motor_current_ma max_count:4

# Maximum OTA chunk size (1KB per design doc)
OtaWriteChunkCmd.data max_size:1024

# Enable fixed-length encoding
* int_size:IS_32
\end{lstlisting}

%=============================================================================
\section{Phase 3: Frame Layer}
%=============================================================================

\subsection{Frame Structure}

\begin{figure}[h]
\centering
\begin{tikzpicture}[
    box/.style={draw, rectangle, minimum height=0.8cm, font=\small},
    node distance=0cm
]
    \node[box, minimum width=1.5cm, fill=blue!20] (sync) {SYNC};
    \node[box, minimum width=1.2cm, right=of sync, fill=green!20] (seq) {SEQ};
    \node[box, minimum width=1.2cm, right=of seq, fill=green!20] (len) {LEN};
    \node[box, minimum width=1.0cm, right=of len, fill=green!20] (type) {TYPE};
    \node[box, minimum width=1.0cm, right=of type, fill=green!20] (flags) {FLAGS};
    \node[box, minimum width=3.0cm, right=of flags, fill=yellow!20] (payload) {PROTOBUF};
    \node[box, minimum width=1.5cm, right=of payload, fill=red!20] (crc) {CRC-32};

    % Dimensions
    \draw[<->] ([yshift=-0.5cm]sync.south west) -- node[below] {2B} ([yshift=-0.5cm]sync.south east);
    \draw[<->] ([yshift=-0.5cm]seq.south west) -- node[below] {2B} ([yshift=-0.5cm]seq.south east);
    \draw[<->] ([yshift=-0.5cm]len.south west) -- node[below] {2B} ([yshift=-0.5cm]len.south east);
    \draw[<->] ([yshift=-0.5cm]type.south west) -- node[below] {1B} ([yshift=-0.5cm]type.south east);
    \draw[<->] ([yshift=-0.5cm]flags.south west) -- node[below] {1B} ([yshift=-0.5cm]flags.south east);
    \draw[<->] ([yshift=-0.5cm]payload.south west) -- node[below] {variable} ([yshift=-0.5cm]payload.south east);
    \draw[<->] ([yshift=-0.5cm]crc.south west) -- node[below] {4B} ([yshift=-0.5cm]crc.south east);
\end{tikzpicture}
\caption{Frame Structure (Total: 12 + payload bytes)}
\end{figure}

\subsection{Frame Header Fields}

\begin{table}[h]
\centering
\caption{Frame Header Specification}
\begin{tabular}{llp{6cm}}
\toprule
\textbf{Field} & \textbf{Size} & \textbf{Description} \\
\midrule
SYNC & 2 bytes & 0xAA55 (frame synchronization marker) \\
SEQ & 2 bytes & Sequence number (0-65535, wraps) \\
LEN & 2 bytes & Payload length (0-1024 for OTA chunks) \\
TYPE & 1 byte & Message type (0x01=CMD, 0x02=RESP, 0x03=ACK, 0x04=NACK) \\
FLAGS & 1 byte & Control flags (reserved for future use) \\
PAYLOAD & variable & nanopb-encoded protobuf message \\
CRC-32 & 4 bytes & IEEE 802.3 checksum (hardware-accelerated) \\
\bottomrule
\end{tabular}
\end{table}

\subsection{Frame Layer API}

\begin{lstlisting}[caption={star\_frame.h - Frame Layer Interface}]
#ifndef STAR_FRAME_H
#define STAR_FRAME_H

#include <stdint.h>
#include <stddef.h>
#include <stdbool.h>
#include "esp_err.h"

#define STAR_FRAME_SYNC     0xAA55
#define STAR_FRAME_HDR_SIZE 8       // SYNC(2) + SEQ(2) + LEN(2) + TYPE(1) + FLAGS(1)
#define STAR_FRAME_CRC_SIZE 4
#define STAR_FRAME_MAX_PAYLOAD 1040 // 1KB OTA chunk + margin

typedef enum {
    STAR_FRAME_TYPE_CMD  = 0x01,
    STAR_FRAME_TYPE_RESP = 0x02,
    STAR_FRAME_TYPE_ACK  = 0x03,
    STAR_FRAME_TYPE_NACK = 0x04,
} star_frame_type_t;

/**
 * @brief Build frame from payload
 * @param payload Protobuf-encoded payload
 * @param payload_len Length of payload
 * @param seq Sequence number
 * @param type Frame type
 * @param out_frame Output frame buffer
 * @param out_len Output frame length
 * @return ESP_OK on success
 */
esp_err_t star_frame_build(const uint8_t* payload, size_t payload_len,
                           uint16_t seq, uint8_t type,
                           uint8_t* out_frame, size_t* out_len);

/**
 * @brief Parse frame and extract payload
 * @param frame Input frame buffer
 * @param frame_len Length of frame
 * @param out_payload Output payload buffer
 * @param out_payload_len Output payload length
 * @param out_seq Output sequence number
 * @param out_type Output frame type
 * @return ESP_OK on success
 */
esp_err_t star_frame_parse(const uint8_t* frame, size_t frame_len,
                           uint8_t* out_payload, size_t* out_payload_len,
                           uint16_t* out_seq, uint8_t* out_type);

/**
 * @brief Validate frame CRC-32
 * @param frame Frame buffer
 * @param frame_len Frame length
 * @return true if CRC valid, false otherwise
 */
bool star_frame_validate(const uint8_t* frame, size_t frame_len);

#endif // STAR_FRAME_H
\end{lstlisting}

%=============================================================================
\section{Phase 4: ARQ Layer}
%=============================================================================

\subsection{Stop-and-Wait State Machine}

\begin{figure}[h]
\centering
\begin{tikzpicture}[
    state/.style={circle, draw, minimum size=1.5cm, font=\small},
    auto,
    node distance=3cm,
    >=latex
]
    \node[state] (idle) {IDLE};
    \node[state, right=of idle] (send) {SEND};
    \node[state, right=of send] (wait) {WAIT\_ACK};
    \node[state, below=of wait] (retry) {RETRY};

    \draw[->] (idle) -- node {packet ready} (send);
    \draw[->] (send) -- node {transmitted} (wait);
    \draw[->] (wait) -- node[above] {ACK OK} (idle);
    \draw[->] (wait) -- node {timeout} (retry);
    \draw[->] (wait) -- node[left] {NACK} (retry);
    \draw[->] (retry) -- node[below] {retry < MAX} (send);
    \draw[->] (retry.west) -- node[left] {retry = MAX} ++(-2,0) |- (idle);
\end{tikzpicture}
\caption{Stop-and-Wait Sender State Machine}
\end{figure}

\subsection{ARQ Configuration}

\begin{table}[h]
\centering
\caption{ARQ Parameters}
\begin{tabular}{lll}
\toprule
\textbf{Parameter} & \textbf{Value} & \textbf{Rationale} \\
\midrule
Timeout & 8~ms & Design doc line 835 \\
Max retries & 3 & Balance reliability vs latency \\
Sequence range & 0-65535 & 16-bit wraparound \\
Window size & 1 & Stop-and-Wait (27× margin) \\
\bottomrule
\end{tabular}
\end{table}

\subsection{ARQ Layer API}

\begin{lstlisting}[caption={star\_arq.h - ARQ Layer Interface}]
#ifndef STAR_ARQ_H
#define STAR_ARQ_H

#include <stdint.h>
#include <stddef.h>
#include "esp_err.h"
#include "star_frame.h"

#define STAR_ARQ_TIMEOUT_MS 8
#define STAR_ARQ_MAX_RETRIES 3

typedef struct star_arq_handle_s star_arq_handle_t;

typedef enum {
    STAR_ARQ_STATE_IDLE,
    STAR_ARQ_STATE_SEND,
    STAR_ARQ_STATE_WAIT_ACK,
    STAR_ARQ_STATE_RETRY,
} star_arq_state_t;

/**
 * @brief Initialize ARQ layer
 * @param handle ARQ handle
 * @return ESP_OK on success
 */
esp_err_t star_arq_init(star_arq_handle_t* handle);

/**
 * @brief Send payload with ARQ
 * @param handle ARQ handle
 * @param payload Protobuf-encoded payload
 * @param len Payload length
 * @param type Frame type
 * @return ESP_OK on success, ESP_FAIL on max retries
 */
esp_err_t star_arq_send(star_arq_handle_t* handle,
                        const uint8_t* payload, size_t len,
                        uint8_t type);

/**
 * @brief Receive payload with ARQ
 * @param handle ARQ handle
 * @param payload Output payload buffer
 * @param len Output payload length
 * @param type Output frame type
 * @param timeout_ms Receive timeout
 * @return ESP_OK on success
 */
esp_err_t star_arq_receive(star_arq_handle_t* handle,
                           uint8_t* payload, size_t* len,
                           uint8_t* type, uint32_t timeout_ms);

/**
 * @brief Process ACK/NACK
 * @param handle ARQ handle
 * @param seq Sequence number from ACK/NACK
 * @param is_ack true for ACK, false for NACK
 * @return ESP_OK on success
 */
esp_err_t star_arq_process_ack(star_arq_handle_t* handle,
                               uint16_t seq, bool is_ack);

/**
 * @brief Get current sequence number
 * @param handle ARQ handle
 * @return Current sequence number
 */
uint16_t star_arq_get_seq(star_arq_handle_t* handle);

#endif // STAR_ARQ_H
\end{lstlisting}

%=============================================================================
\section{Phase 5: Communication Task}
%=============================================================================

\subsection{Architecture}

\begin{figure}[h]
\centering
\begin{tikzpicture}[
    box/.style={draw, rectangle, minimum width=6cm, minimum height=1cm, align=center},
    arrow/.style={->, >=latex, thick},
    node distance=1.2cm
]
    \node[box, fill=blue!20] (task) {communication\_task()};
    \node[box, fill=green!20, below=of task] (arq_rx) {star\_arq\_receive()};
    \node[box, fill=yellow!20, below=of arq_rx] (frame_parse) {star\_frame\_parse()};
    \node[box, fill=orange!20, below=of frame_parse] (pb_decode) {pb\_decode()};
    \node[box, fill=red!20, below=of pb_decode] (handler) {Handle command};
    \node[box, fill=orange!20, below=of handler] (pb_encode) {pb\_encode()};
    \node[box, fill=yellow!20, below=of pb_encode] (frame_build) {star\_frame\_build()};
    \node[box, fill=green!20, below=of frame_build] (arq_tx) {star\_arq\_send()};

    \draw[arrow] (task) -- (arq_rx);
    \draw[arrow] (arq_rx) -- (frame_parse);
    \draw[arrow] (frame_parse) -- (pb_decode);
    \draw[arrow] (pb_decode) -- (handler);
    \draw[arrow] (handler) -- (pb_encode);
    \draw[arrow] (pb_encode) -- (frame_build);
    \draw[arrow] (frame_build) -- (arq_tx);
\end{tikzpicture}
\caption{Communication Task Flow}
\end{figure}

\subsection{Message Handlers}

\begin{table}[h]
\centering
\caption{Command Handlers}
\begin{tabular}{lll}
\toprule
\textbf{Message} & \textbf{Handler} & \textbf{Response} \\
\midrule
MotorSpeedCmd & queue\_motor\_command() & AckResp \\
MotorStopCmd & queue\_motor\_stop() & AckResp \\
EmergencyStopCmd & immediate\_motor\_stop() & AckResp \\
SetPidGainsCmd & update\_pid\_gains() & AckResp \\
OtaWriteChunkCmd & star\_ota\_write() & AckResp/NackResp \\
RequestTelemetry & encode\_telemetry() & TelemetryResp \\
\bottomrule
\end{tabular}
\end{table}

%=============================================================================
\section{Phase 6: Build Integration}
%=============================================================================

\subsection{platformio.ini Updates}

\begin{lstlisting}[caption={platformio.ini - Add nanopb Dependency}]
[env:esp32s3]
platform = espressif32
board = esp32-s3-devkitc-1
framework = espidf

lib_deps =
    nanopb/Nanopb@^0.4.8

build_flags =
    -DCONFIG_STAR_PROTO_ENABLE=1

extra_scripts =
    pre:scripts/generate_proto.py
\end{lstlisting}

\subsection{Proto Generation Script}

\begin{lstlisting}[caption={scripts/generate\_proto.sh},language=bash]
#!/bin/bash
set -e

PROTO_DIR="lib/star_proto/proto"
OUT_DIR="lib/star_proto/src/generated"

# Find nanopb generator
NANOPB_GENERATOR=$(find ~/.platformio/packages -name "protoc-gen-nanopb" | head -n 1)

if [ -z "$NANOPB_GENERATOR" ]; then
    echo "Error: nanopb generator not found"
    exit 1
fi

# Generate C code from .proto
mkdir -p "$OUT_DIR"
protoc --plugin=protoc-gen-nanopb="$NANOPB_GENERATOR" \
       --nanopb_out="$OUT_DIR" \
       -I"$PROTO_DIR" \
       "$PROTO_DIR/messages.proto"

echo "Proto generation complete: $OUT_DIR"
\end{lstlisting}

%=============================================================================
\section{Implementation Checklist}
%=============================================================================

\begin{enumerate}[label=\protect\Large$\square$]
    \item Phase 1: Create lib/star\_crc/ with hardware CRC-32
    \item Phase 2: Create lib/star\_proto/ directory structure
    \item Phase 3: Write messages.proto and messages.options
    \item Phase 4: Implement star\_frame.c (frame layer)
    \item Phase 5: Implement star\_arq.c (ARQ layer)
    \item Phase 6: Rewrite communication\_task.c
    \item Phase 7: Update platformio.ini with nanopb
    \item Phase 8: Create scripts/generate\_proto.sh
    \item Phase 9: Integration testing with RPi5
    \item Phase 10: Performance validation (100~Hz, <20~ms emergency stop)
\end{enumerate}

%=============================================================================
\section{Success Criteria}
%=============================================================================

\begin{table}[h]
\centering
\caption{Validation Metrics}
\begin{tabular}{lcc}
\toprule
\textbf{Metric} & \textbf{Target} & \textbf{Measurement Method} \\
\midrule
Memory usage & <10~KB & ESP-IDF heap/stack profiling \\
Reliability & 99.9999\% & CRC-32 error injection test \\
Throughput & 100~Hz & Sustained motor command rate \\
Emergency stop & <20~ms & Logic analyzer timing \\
ARQ retry & >95\% success & Error injection test \\
\bottomrule
\end{tabular}
\end{table}

%=============================================================================
\section{References}
%=============================================================================

\begin{enumerate}
    \item \texttt{Protobuf\_Protocol\_Design\_Analysis.tex} - System design document
    \item nanopb Documentation - https://jpa.kapsi.fi/nanopb/
    \item ESP32-S3 Technical Reference Manual - Espressif Systems
    \item Protocol Buffers Language Guide - Google
\end{enumerate}

